\section{Resultados}
\label{sec:resultados}

\subsection{Resultados para el esquema de cuantización uniforme}
\label{sec:uniform_results}

En primer lugar, en la \reffig{uniform_results} se presentan los resultados obtenidos con el esquema de cuantización uniforme, donde se grafica la precisión alcanzada en función de la complejidad del modelo cuantizado.

\defaultfigure{resultados/complexity_vs_quantized_uniform_stats.png}{Resultados de la cuantización uniforme.}{uniform_results}

En la tabla \ref{tab:uniform_results} se mostraron los resultados cuantitativos obtenidos con el esquema de cuantización uniforme.

{\footnotesize
\begin{table}[H]
    \centering
    \caption{Resultados de cuantización uniforme}
    \begin{tabular}{@{}r r r r c@{}}
        \toprule
        Prec. (máx.) & Prec. Rel. (\%) & Compl. (Kb) & Compres. & Bits \\
        \midrule
        $0.7591$ & $102.51$ & $50917.08$ & $1.34\times$ & $24$ \\
        $0.7585$ & $102.43$ & $16997.08$ & $4.00\times$ & $8$ \\
        $0.7584$ & $102.41$ & $33957.08$ & $2.00\times$ & $16$ \\
        $0.7581$ & $102.38$ & $21237.08$ & $3.20\times$ & $10$ \\
        $0.7580$ & $102.36$ & $12757.08$ & $5.33\times$ & $6$ \\
        $0.7536$ & $101.77$ & $10637.08$ & $6.39\times$ & $5$ \\
        \rowcolor{gray!20}
        $0.7405$ & $100.00$ & $67988.31$ & $1.00\times$ & \texttt{-} \\
        $0.7367$ & $99.49$ & $8517.08$ & $7.98\times$ & $4$ \\
        $0.6399$ & $86.38$ & $4277.08$ & $15.90\times$ & $2$ \\
        $0.6371$ & $86.03$ & $6397.08$ & $10.63\times$ & $3$ \\
        $0.3698$ & $49.94$ & $2157.08$ & $31.52\times$ & $1$ \\
        \bottomrule
    \end{tabular}
    \label{tab:uniform_results}
\end{table}
}


Los resultados evidencian una buena capacidad de compresión hasta aproximadamente $8\times$, con una precisión cercana al $99\%$ del modelo original. Por ejemplo, con $8$ bits se alcanzó una precisión máxima de $0.7585$ ($4.00\times$). Con $10$ y $16$ bits, con valores máximos de $0.7581$ y $0.7584$, con compresiones de $3.20\times$ y $2.00\times$, respectivamente.

A medida que disminuye el ancho de bits, la precisión disminuye de manera significativa (al utilizar $3$ bits, la precisión descendió a $86.03\%$ en comparación con el modelo original). Esto evidencia que la cuantización uniforme no es adecuada para realizar compresiones agresivas.

Asimismo, hubo configuraciones que superaron la precisión del modelo original. Al continuar el entrenamiento tras la cuantización y con resolución relativamente alta ($\geq 8$ bits), el modelo se optimiza y ajusta los pesos con más datos y alcanza una precisión ligeramente superior.

\subsection{Resultados para el esquema de cuantización flexible}
\label{sec:flexible_results}

\defaultfigure{resultados/complexity_vs_quantized_flex_stats2.png}{Resultados de la cuantización flexible.}{complexity_vs_quantized_flex_stats2}

A continuación, la tabla \ref{tab:flexible_results} presenta los resultados detallados del esquema de cuantización flexible, que muestra diferentes combinaciones de bits y niveles de cuantización. Este esquema permite un control más fino del balance entre precisión y complejidad, lo que habilita configuraciones con alta compresión y pérdidas de precisión contenidas.


{\footnotesize
\begin{longtable}{@{}r r r r c c@{}}
    \caption{Resultados de cuantización flexible} \label{tab:flexible_results} \\
    \toprule
    Precisión & Prec. Rel. (\%) & Compl. (Kb) & Compres. & Bits & Niv. \\
    \midrule
    \endfirsthead

    \multicolumn{6}{c}{\tablename\ \thetable\ -- Continuación de la página anterior} \\
    \toprule
    Precisión & Prec. Rel. (\%) & Compl. (Kb) & Compres. & Bits & Niv. \\
    \midrule
    \endhead

    \midrule
    \multicolumn{6}{r@{}}{Continúa en la siguiente página...} \\
    \endfoot

    \bottomrule
    \endlastfoot

    \rowcolor{gray!20}
    $0.7437$ & $100.00$ & $67988.31$ & $1.00\times$ & \texttt{-} & \texttt{-} \\
    $0.7354$ & $98.88$ & $7633.31$ & $8.91\times$ & $6$ & $25$ \\
    $0.7342$ & $98.72$ & $5369.26$ & $12.66\times$ & $6$ & $7$ \\
    $0.7338$ & $98.67$ & $8596.11$ & $7.91\times$ & $8$ & $30$ \\
    $0.7328$ & $98.53$ & $6598.73$ & $10.30\times$ & $6$ & $16$ \\
    $0.7298$ & $98.13$ & $5990.61$ & $11.35\times$ & $6$ & $10$ \\
    $0.7294$ & $98.08$ & $7120.25$ & $9.55\times$ & $6$ & $20$ \\
    $0.7290$ & $98.02$ & $4279.46$ & $15.89\times$ & $4$ & $7$ \\
    $0.7289$ & $98.01$ & $4941.02$ & $13.76\times$ & $4$ & $10$ \\
    $0.7287$ & $97.98$ & $7245.11$ & $9.38\times$ & $8$ & $16$ \\
    $0.7286$ & $97.97$ & $7891.29$ & $8.62\times$ & $6$ & $30$ \\
    $0.7284$ & $97.94$ & $5099.06$ & $13.33\times$ & $4$ & $8$ \\
    $0.7284$ & $97.94$ & $8409.12$ & $8.09\times$ & $8$ & $25$ \\
    $0.7278$ & $97.86$ & $7811.66$ & $8.70\times$ & $8$ & $20$ \\
    $0.7270$ & $97.75$ & $4845.85$ & $14.03\times$ & $4$ & $6$ \\
    $0.7267$ & $97.71$ & $5539.91$ & $12.27\times$ & $6$ & $8$ \\
    $0.7264$ & $97.67$ & $5380.73$ & $12.64\times$ & $8$ & $7$ \\
    $0.7253$ & $97.53$ & $5902.68$ & $11.52\times$ & $8$ & $8$ \\
    $0.7245$ & $97.42$ & $3884.17$ & $17.50\times$ & $8$ & $4$ \\
    $0.7244$ & $97.40$ & $6533.89$ & $10.41\times$ & $8$ & $10$ \\
    $0.7240$ & $97.35$ & $4897.42$ & $13.88\times$ & $8$ & $6$ \\
    $0.7223$ & $97.12$ & $4187.65$ & $16.24\times$ & $8$ & $5$ \\
    $0.7223$ & $97.12$ & $4955.85$ & $13.72\times$ & $6$ & $6$ \\
    $0.7222$ & $97.11$ & $3650.13$ & $18.63\times$ & $4$ & $4$ \\
    $0.7211$ & $96.96$ & $4139.22$ & $16.43\times$ & $6$ & $5$ \\
    $0.7209$ & $96.93$ & $2953.72$ & $23.02\times$ & $4$ & $3$ \\
    $0.7203$ & $96.85$ & $4041.94$ & $16.82\times$ & $4$ & $5$ \\
    $0.7196$ & $96.76$ & $3362.26$ & $20.22\times$ & $6$ & $4$ \\
    $0.7194$ & $96.73$ & $2805.85$ & $24.23\times$ & $8$ & $3$ \\
    $0.7178$ & $96.52$ & $2677.41$ & $25.39\times$ & $6$ & $3$ \\
    $0.7061$ & $94.94$ & $2389.88$ & $28.45\times$ & $8$ & $2$ \\
    $0.7054$ & $94.85$ & $2484.54$ & $27.36\times$ & $6$ & $2$ \\
    $0.6987$ & $93.95$ & $2119.35$ & $32.08\times$ & $4$ & $2$ \\
\end{longtable}
}


A compresiones muy altas (por ejemplo, $8/3$ con $24.23\times$), la precisión se mantuvo alrededor de $0.7194$, lo que muestra que este esquema con cuantización flexible es capaz de sostener un desempeño razonable incluso bajo compresiones agresivas.

\subsection{Comparación de esquemas de cuantización}

Finalmente, en la \reffig{uniform_vs_flex_max_curves} se muestran los resultados obtenidos con los dos esquemas de cuantización utilizados, tanto el esquema de cuantización uniforme como el esquema de cuantización flexible.

\defaultfigure{resultados/uniform_vs_flex_max_curves.png}{Comparación de los esquemas de cuantización uniforme y flexible.}{uniform_vs_flex_max_curves}

En la zona de mayor tasa de compresión, se observa que el esquema de cuantización flexible provee mejores valores de precisión en comparación con el esquema de cuantización uniforme. La cuantización uniforme no permite altas tasas de compresión sin una pérdida significativa de precisión, como se adelant\'o en la \ref{sec:estado_arte}.

En la tabla \ref{tab:compare_results} se exponen los resultados cuantitativos completos que mostrados en la gráfica de la \reffig{uniform_vs_flex_max_curves}, incluyendo todos los modelos entrenados con ambos esquemas de cuantización. En términos de la frontera de Pareto (precisión vs. complejidad), el esquema flexible domina a partir de compresiones superiores a \mbox{$\sim$10$\times$}. También se observan configuraciones con 12--18$\times$ y pérdidas de precisión por debajo de 2 puntos porcentuales respecto del modelo original. Para compresiones moderadas (\mbox{$\leq$ 8$\times$}), ambos esquemas muestran desempeños similares, mientras que para tasas de compresión muy altas (\mbox{$\geq$ 20$\times$}) la ventaja del esquema flexible es aún más pronunciada.

{\footnotesize
\begin{longtable}{@{}r r r r c c l@{}}
    \caption{Comparación de precisión y complejidad para diferentes esquemas de cuantización} \label{tab:compare_results} \\
    \toprule
    Precisión & Prec. Rel. (\%) & Compl. (Kb) & Compres. & Bits & Niv. & Tipo \\
    \midrule
    \endfirsthead

    \multicolumn{7}{c}{\tablename\ \thetable\ -- Continuación de la página anterior} \\
    \toprule
    Precisión & Prec. Rel. (\%) & Compl. (Kb) & Compres. & Bits & Niv. & Tipo \\
    \midrule
    \endhead

    \midrule
    \multicolumn{7}{r@{}}{Continúa en la siguiente página...} \\
    \endfoot

    \bottomrule
    \endlastfoot

    $0.7591$ & $102.07$ & $50917.08$ & $1.34\times$ & $24$ & \texttt{-} & uniforme \\
    $0.7585$ & $101.99$ & $16997.08$ & $4.00\times$ & $8$ & \texttt{-} & uniforme \\
    $0.7584$ & $101.98$ & $33957.08$ & $2.00\times$ & $16$ & \texttt{-} & uniforme \\
    $0.7581$ & $101.94$ & $21237.08$ & $3.20\times$ & $10$ & \texttt{-} & uniforme \\
    $0.7580$ & $101.92$ & $12757.08$ & $5.33\times$ & $6$ & \texttt{-} & uniforme \\
    $0.7536$ & $101.33$ & $10637.08$ & $6.39\times$ & $5$ & \texttt{-} & uniforme \\
    \rowcolor{gray!20}
    $0.7437$ & $100.00$ & $67988.31$ & $1.00\times$ & \texttt{-} & \texttt{-} & original \\
    $0.7367$ & $99.06$ & $8517.08$ & $7.98\times$ & $4$ & \texttt{-} & uniforme \\
    $0.7354$ & $98.88$ & $7633.31$ & $8.91\times$ & $6$ & $25$ & flexible \\
    $0.7342$ & $98.72$ & $5369.26$ & $12.66\times$ & $6$ & $7$ & flexible \\
    $0.7338$ & $98.67$ & $8596.11$ & $7.91\times$ & $8$ & $30$ & flexible \\
    $0.7328$ & $98.53$ & $6598.73$ & $10.30\times$ & $6$ & $16$ & flexible \\
    $0.7298$ & $98.13$ & $5990.61$ & $11.35\times$ & $6$ & $10$ & flexible \\
    $0.7294$ & $98.08$ & $7120.25$ & $9.55\times$ & $6$ & $20$ & flexible \\
    $0.7290$ & $98.02$ & $4279.46$ & $15.89\times$ & $4$ & $7$ & flexible \\
    $0.7289$ & $98.01$ & $4941.02$ & $13.76\times$ & $4$ & $10$ & flexible \\
    $0.7287$ & $97.98$ & $7245.11$ & $9.38\times$ & $8$ & $16$ & flexible \\
    $0.7286$ & $97.97$ & $7891.29$ & $8.62\times$ & $6$ & $30$ & flexible \\
    $0.7284$ & $97.94$ & $5099.06$ & $13.33\times$ & $4$ & $8$ & flexible \\
    $0.7284$ & $97.94$ & $8409.12$ & $8.09\times$ & $8$ & $25$ & flexible \\
    $0.7278$ & $97.86$ & $7811.66$ & $8.70\times$ & $8$ & $20$ & flexible \\
    $0.7270$ & $97.75$ & $4845.85$ & $14.03\times$ & $4$ & $6$ & flexible \\
    $0.7267$ & $97.71$ & $5539.91$ & $12.27\times$ & $6$ & $8$ & flexible \\
    $0.7264$ & $97.67$ & $5380.73$ & $12.64\times$ & $8$ & $7$ & flexible \\
    $0.7253$ & $97.53$ & $5902.68$ & $11.52\times$ & $8$ & $8$ & flexible \\
    $0.7245$ & $97.42$ & $3884.17$ & $17.50\times$ & $8$ & $4$ & flexible \\
    $0.7244$ & $97.40$ & $6533.89$ & $10.41\times$ & $8$ & $10$ & flexible \\
    $0.7240$ & $97.35$ & $4897.42$ & $13.88\times$ & $8$ & $6$ & flexible \\
    $0.7223$ & $97.12$ & $4187.65$ & $16.24\times$ & $8$ & $5$ & flexible \\
    $0.7223$ & $97.12$ & $4955.85$ & $13.72\times$ & $6$ & $6$ & flexible \\
    $0.7222$ & $97.11$ & $3650.13$ & $18.63\times$ & $4$ & $4$ & flexible \\
    $0.7211$ & $96.96$ & $4139.22$ & $16.43\times$ & $6$ & $5$ & flexible \\
    $0.7209$ & $96.93$ & $2953.72$ & $23.02\times$ & $4$ & $3$ & flexible \\
    $0.7203$ & $96.85$ & $4041.94$ & $16.82\times$ & $4$ & $5$ & flexible \\
    $0.7196$ & $96.76$ & $3362.26$ & $20.22\times$ & $6$ & $4$ & flexible \\
    $0.7194$ & $96.73$ & $2805.85$ & $24.23\times$ & $8$ & $3$ & flexible \\
    $0.7178$ & $96.52$ & $2677.41$ & $25.39\times$ & $6$ & $3$ & flexible \\
    $0.7061$ & $94.94$ & $2389.88$ & $28.45\times$ & $8$ & $2$ & flexible \\
    $0.7054$ & $94.85$ & $2484.54$ & $27.36\times$ & $6$ & $2$ & flexible \\
    $0.6987$ & $93.95$ & $2119.35$ & $32.08\times$ & $4$ & $2$ & flexible \\
    $0.6399$ & $86.04$ & $4277.08$ & $15.90\times$ & $2$ & \texttt{-} & uniforme \\
    $0.6371$ & $85.67$ & $6397.08$ & $10.63\times$ & $3$ & \texttt{-} & uniforme \\
    $0.3698$ & $49.72$ & $2157.08$ & $31.52\times$ & $1$ & \texttt{-} & uniforme \\
\end{longtable}
}

